\subsection{Анализ предметной области и формирование литературной базы}

На начальном этапе работы был проведен систематический поиск и анализ научно-технической литературы по тематике датчиков угловых перемещений, применяемых в высокоточных оптических и астрономических системах. Изучены принципы построения абсолютных и инкрементных энкодеров, рассмотрены основные физические методы измерения угла — оптический, магнитный, индуктивный и емкостный — а также их сравнительные характеристики применительно к задачам точного позиционирования. На основании проведенного анализа сформирована теоретическая база работы и составлен литературный обзор, включенный в настоящую работу в качестве отдельного раздела.

\subsection{Уточнение постановки задачи}
В рамках научно‑технического сотрудничества был визит в АО «Мостком» (г. Рязань) — компанию, специализирующуюся на разработке прецизионных встраиваемых компонентов и систем управления. Во время обсуждения были получены дополнительные материалы по тематике исследования, предоставлен исходный код действующей системы считывания и обработки сигнала энкодера, а также показаны реальная конструкция датчика и узел его установки.

После совместного обсуждения и анализа этих материалов исходная постановка задачи была уточнена. Основное внимание было сосредоточено на алгоритмах самокомпенсации и самокоррекции систематических ошибок энкодера с учетом условий работы в составе астрономических телескопов. К ключевым ограничениям относятся высокие требования к точности позиционирования (от единиц до десятых долей угловой секунды), невозможность прерывать наблюдения для повторной калибровки и ограниченный диапазон поворота оси телескопа, из‑за чего нельзя напрямую применять алгоритмы, рассчитанные на полнооборотные и многооборотные системы.

\subsection{Специфика параметров прототипа}

При разработке прототипа энкодера для телескопа АЗТ-6 в данной работе действует ряд специфических ограничений.

\begin{enumerate}
    \item \textbf{Метод изготовления растра}\label{point_raster}. Характерная точность позиционирования экструдера 3D-принтера составляет $0.1$–$0.2$ мм --- это дает ограничение на минимальный размер растра.
    \item \textbf{Компактность конструкции, расположения датчика на осях монтировки АЗТ-6}. Максимальные наружные габариты растра существующими посадочными диаметрами и монтажными зазорами телескопа, что также ограничивает радиус шкалы.
    \item \textbf{Конструкция считывающей головки энкодера}\label{point_lense_distortion}. С учетом \hyperref[point_raster]{п.~1} было принято решение использовать линзу для проецирования растра на ПЗС-матрицу, что позволяет эффективно увеличить используемую длину матрицы ценой внесения оптических искажений. Основным требованием к оптической схеме является обеспечение достаточного числа пикселей ПЗС на один бит растра, необходимого для последующего анализа сигнала.
\end{enumerate}

По результатам тестовых печатей и измерений были выбраны следующие масштабы:
\begin{table}[ht!]
    \centering
    \begin{tabular}{c|c|c}
        \hline
        Параметр & Оценка & Значение \\
        \hline
        $k$, пикс. & $\geq 3$ & 6\\
        $\sigma(k)$, пикс. & $\leq 0.1$ & -\\
        $L$, мм & $\leq 30$& 24\\
        $R$, мм & $\leq 100$ & 100\\
        $\Delta R$, мм & $\geq5$ & 5\\
        $w$, мм & $\geq 0.5$ & 1.6\\
        $\lambda = \frac{k}{w}$, $\frac{\text{пикс}}{\text{мм}}$ & $\geq 3$ & 3.75\\
    \end{tabular}
    \caption{Оцененные параметры установки}
    \label{tab:work_values}
\end{table}

В соответствии с условиями \eqref{eq:resolution}, \eqref{eq:code_len}, \eqref{eq:estimate_dR}, выполним проверку корректности и оценим характеристики:
\begin{equation*}
    \rho \simeq \frac{0.1}{100\times3.75}~\text{рад} \simeq 55'' 
\end{equation*}

Положим, что допустимое отклонение матрицы от нормального положения не превышает $\alpha = 5^\circ$. Получим ограничение на $\Delta R$:
\begin{equation*}
    \Delta R \gtrsim 3.57~\text{мм}
\end{equation*}

Для выбранного \ref{tab:work_values} $\Delta R = 5$ мм условие выполняется. Оценим число бит в заданной конфигурации:
\begin{equation*}
    n \geq \log_2\left(27.49\cdot(n+1)\right) \Rightarrow n \geq 7.94
\end{equation*}

При заданных $k$, $L$, $w$, $R$ в поле зрения матрицы (в пределах зоны умеренной дисторсии и виньетирования) умещается 15 бит кода. Таким образом, полное число кодов составляет $N = 27.49 \times 15 \simeq 412$.

Поскольку $2^8 < 412 < 2^9$, в качестве используемого значения было выбрано $n = 9$.
\subsection{Проектирование прототипа и анализ технических ограничений}

\subsubsection{Моделирование установки энкодера на телескопе АЗТ-6}

В качестве объекта применения разрабатываемого прототипа выбран астрономический телескоп АЗТ-6. На данном этапе было выполнено моделирование (см. \hyperref[sec:application2]{Приложение 2}) геометрии расположения датчика угла на полярной оси монтировки телескопа с целью определения допустимых габаритов устройства, требований к допустимым положениям считывающих головок и конструктивных ограничений, обусловленных существующей механической компоновкой инструмента.

\subsubsection{Подбор компонентов и анализ ограничений прототипа}

На основании результатов моделирования и технических требований был проведен анализ возможных компонентов для изготовления прототипа. Принципиальным ограничением являлось то, что измерительный растр предполагалось изготавливать методом 3D-печати. Это позволяет значительно сократить сроки изготовления для тестирования, однако накладывает ряд ограничений: точность геометрии ограничена возможностями принтера, крупные элементы подвержены деформации при охлаждении.

В рамках указанных ограничений выполнены тестовые печати линейных растров с периодической структурой (чередование прозрачных и непрозрачных участков) с последующей оценкой качества. По результатам испытаний (см. \hyperref[sec:application1]{Приложение 1}) определен масштаб кодового растра, воспроизводимый методом 3D-печати.

\subsection{Модель дисторсии и ее коррекция} \label{subsec:distortion_model}

\subsubsection{Обоснование необходимости коррекции}

Как отмечено в \hyperref[point_lense_distortion]{п.~3} специфики прототипа, использование линзы вносит радиальную дисторсию --- геометрическое искажение, при котором масштаб изображения зависит от расстояния до оптической оси. В результате равномерная структура растра проецируется на матрицу с переменным масштабом, что приводит к систематической ошибке при определении положений фронтов и, как следствие, к ошибке в восстановлении углового положения.

Для улучшения качества была сделана оценка параметров дисторсии и коррекция структуры считываемого сигнала до этапа восстановления битовой последовательности.

\subsubsection{Математическая модель}

В данной работе используется модель деления (division model) дисторсии, описанная в работе \cite{aleman_lens_distortion}. В одномерном случае, то есть вдоль линейки ПЗС-матрицы, модель записывается в виде:
\begin{equation}
x_d = \frac{x_u}{1 + k\,x_u^2},
\label{eq:distortion_forward}
\end{equation}
где $x_u$ --- неискаженная координата точки изображения, $x_d$ --- наблюдаемая координата на сенсоре, $k$ --- коэффициент дисторсии.

Для обеспечения независимости коэффициента $k$ от числа пикселей линейки координаты переводятся из пиксельного пространства в нормированное на $[-1,~1]$ относительно центра кадра.

При $N \geq 2$ крайние пиксели сенсора соответствуют нормированным координатам, близким к $x=\pm1$. Такая нормировка позволяет использовать единый коэффициент дисторсии для линеек различной длины.

\subsubsection{Обратное преобразование}

Для восстановления неискаженных координат по наблюдаемым необходимо обратить преобразование \eqref{eq:distortion_forward}. Подстановка $x_d(x_u)$ приводит к квадратному уравнению относительно $x_u$:
\begin{equation}
k\,x_d\,x_u^2 - x_u + x_d = 0.
\label{eq:distortion_quadratic}
\end{equation}

Физически корректный корень, удовлетворяющий условию $x_u \to x_d$ при $k \to 0$, имеет вид
\begin{equation}
x_u = \frac{1 - \sqrt{1 - 4k\,x_d^2}}{2k\,x_d}.
\label{eq:distortion_inverse_raw}
\end{equation}

Однако данная форма численно неустойчива при малых $k$, поскольку содержит вычитание близких величин в знаменателе. Поэтому в реализации используется эквивалентная устойчивая запись:
\begin{equation}
x_u = \frac{2x_d}{1 + \sqrt{1 - 4k\,x_d^2}}.
\label{eq:distortion_inverse_stable}
\end{equation}

При численной реализации выражение под корнем ограничивается снизу нулем, что исключает появление комплексных значений из-за накопления ошибок округления. Область определения обратного преобразования задается условием неотрицательности дискриминанта:
\begin{equation}
    1 - 4k\,x_d^2 \geq 0 \quad \Longrightarrow \quad |x_d| \leq \frac{1}{2\sqrt{k}}.
    \label{eq:distortion_domain}
\end{equation}

Вычисление производится для нормированных координат. При $|x_d| \leq 1$ это дает ограничение $k \leq 0.25$.

\subsubsection{Калибровка по образу периодической структуры}

Идея калибровки основана на том, что после устранения дисторсии расстояния между фронтами периодической структуры должны быть одинаковыми. В качестве калибровочного объекта используется последовательность чередующихся полос равной ширины, назовем этот объект \textit{зеброй} (см. рис \ref{fig:zebra}).

Пусть $\{x_{d, \mathrm{px}}^{(i)}\}_{i=1}^{M}$ --- координаты найденных фронтов зебры в пиксельных координатах сенсора, определенные с субпиксельной точностью. Далее они переводятся в нормированное пространство:
\begin{equation}
x_d^{(i)} = \frac{x_{d,\mathrm{px}}^{(i)} - c}{s}.
\label{eq:zebra_edges_norm}
\end{equation}

Для заданного пробного значения $k$ вычисляются соответствующие неискаженные координаты фронтов:
\begin{equation}
\hat{x}_u^{(i)}(k) =
\frac{2x_d^{(i)}}{1 + \sqrt{1 - 4k\left(x_d^{(i)}\right)^2}}.
\label{eq:zebra_undistort}
\end{equation}

После этого они переводятся обратно в пиксельные единицы:
\begin{equation}
\hat{x}_{u,\mathrm{px}}^{(i)}(k) = \hat{x}_u^{(i)}(k)\cdot s + c.
\label{eq:zebra_undistort_px}
\end{equation}

Межфронтовые расстояния:
\begin{equation}
d_i(k) = \hat{x}_{u,\mathrm{px}}^{(i+1)}(k) - \hat{x}_{u,\mathrm{px}}^{(i)}(k),
\quad i=1,\ldots,M-1,
\label{eq:zebra_spacing}
\end{equation}

Коэффициент дисторсии определяется минимизацией нормированной дисперсии межфронтовых расстояний:
\begin{equation}
k^* = \arg\min_{k \in [k_{\min},\,k_{\max}]}
\frac{\mathrm{Var}\,d_i(k)}{\left(\mathrm{E}\,d_i(k)\right)^2}.
\label{eq:distortion_objective}
\end{equation}

Такой функционал не зависит от абсолютного масштаба полос и характеризует только степень неравномерности расстояний между соседними фронтами. В численной реализации поиск минимума выполняется методом ограниченной одномерной оптимизации на отрезке $[k_{\min},\,k_{\max}]$.

\subsubsection{Восстановление сигнала}

После определения коэффициента дисторсии $k^*$ восстановление произвольного сигнала считанного той же матрицей, выполняется как переход от равномерной сетки в искаженном пространстве сенсора к равномерной сетке в неискаженном пространстве изображения.

\paragraph{Две постановки задачи восстановления.}

Задача восстановления может быть сформулирована двумя эквивалентными способами.

\begin{enumerate}
\item \textbf{Прямая постановка.}  
Задается равномерная выходная сетка в неискаженном пространстве $\{x_u^{(j)}\}$, после чего для каждой ее точки по прямой модели
\begin{equation}
x_d^{(j)} = \frac{x_u^{(j)}}{1 + k^*\left(x_u^{(j)}\right)^2}
\label{eq:restore_forward_scheme}
\end{equation}
вычисляется соответствующая координата в искаженном пространстве, и значение сигнала определяется интерполяцией исходного массива в точках $x_d^{(j)}$.

Схема вычислений:
\begin{equation}
x_u^{\mathrm{uniform}} \longrightarrow x_d \longrightarrow\tilde{S}.
\label{eq:restore_scheme_forward}
\end{equation}

\item \textbf{Обратная постановка.}  
Рассматриваются исходные отсчеты сигнала, заданные на равномерной сетке искаженного сенсора. Для каждого исходного отсчета вычисляется его положение в неискаженном пространстве по обратной модели
\begin{equation}
x_u^{(i)} =
\frac{2x_d^{(i)}}{1 + \sqrt{1 - 4k^*\left(x_d^{(i)}\right)^2}},
\label{eq:restore_inverse_scheme}
\end{equation}
после чего сигнал, заданный на полученной неравномерной сетке $\{x_u^{(i)}\}$, интерполируется на новую равномерную сетку в неискаженном пространстве.

Схема вычислений:
\begin{equation}
x_d^{\mathrm{uniform}} \longrightarrow x_u \longrightarrow x_u^{\mathrm{uniform}} \longrightarrow \tilde{S}.
\label{eq:restore_scheme_inverse}
\end{equation}
\end{enumerate}

Обе постановки математически эквивалентны. Однако в данной работе используется именно \textit{обратная постановка}, поскольку исходный сигнал изначально задан на равномерной пиксельной сетке сенсора, и потому естественно сначала вычислить положения этих отсчетов в неискаженном пространстве, а затем выполнить пересчет на новую равномерную сетку.

\paragraph{Алгоритм восстановления, используемый в реализации.}

В реализованном алгоритме восстановление выполняется следующим образом.

\begin{enumerate}
\item Для исходного сигнала задается равномерная сетка отсчетов в пиксельных координатах сенсора:
\begin{equation}
x_{d,\mathrm{px}}^{(i)} = i, \quad i = 0,1,\ldots,N-1.
\label{eq:restore_sensor_grid}
\end{equation}

Соответствующие нормированные координаты определяются соотношением
\begin{equation}
x_d^{(i)} = \frac{x_{d,\mathrm{px}}^{(i)} - c}{s},
\quad
c = \frac{N-1}{2}, \quad s = \max\left(\frac{N-1}{2},\,1\right).
\label{eq:restore_sensor_norm}
\end{equation}

\item Для каждого исходного отсчета вычисляется координата в неискаженном пространстве:
\begin{equation}
x_u^{(i)} =
\frac{2x_d^{(i)}}{1 + \sqrt{1 - 4k^*\left(x_d^{(i)}\right)^2}}.
\label{eq:restore_inverse_main}
\end{equation}

После этого выполняется переход обратно к пиксельным единицам:
\begin{equation}
x_{u,\mathrm{px}}^{(i)} = x_u^{(i)} \cdot s + c.
\label{eq:restore_inverse_px}
\end{equation}

Таким образом, исходный сигнал $S$ начинает рассматриваться как функция, заданная на неравномерной сетке $\{x_{u,\mathrm{px}}^{(i)}\}$.

\item По диапазону найденных координат строится новая равномерная сетка в неискаженном пространстве:
\begin{equation}
x_u^{\min} = \min_i x_{u,\mathrm{px}}^{(i)}, \qquad
x_u^{\max} = \max_i x_{u,\mathrm{px}}^{(i)}.
\label{eq:restore_bounds_current}
\end{equation}

Далее задается равномерная выходная сетка с шагом $\Delta x = \Delta x_{\mathrm{out}}$:
\begin{equation}
x_{u,\mathrm{new}}^{(j)} = x_u^{\min} + j\,\Delta x,
\quad
j = 0,1,\ldots,N'-1,
\label{eq:restore_uniform_grid}
\end{equation}
где
\begin{equation}
N' = \left\lfloor \frac{x_u^{\max} - x_u^{\min}}{\Delta x} \right\rfloor + 1.
\label{eq:restore_size}
\end{equation}

В используемой реализации принимается $\Delta x = 1$ пиксель.

\item Значения восстановленного сигнала получаются интерполяцией исходного массива $S$, заданного на неравномерной сетке $\{x_{u,\mathrm{px}}^{(i)}\}$, в узлах новой равномерной сетки $\{x_{u,\mathrm{new}}^{(j)}\}$:
\begin{equation}
\tilde{s}_j = \mathrm{interp}\left(x_{u,\mathrm{new}}^{(j)}; \left\{x_{u,\mathrm{px}}^{(i)},\,s_i\right\}_{i=0}^{N-1}\right).
\label{eq:restore_interp_current}
\end{equation}
\end{enumerate}

Таким образом, на выходе формируется массив
$\tilde{S} = [\tilde{s}_0, \tilde{s}_1, \ldots, \tilde{s}_{N'-1}]$,
соответствующий равномерной пиксельной сетке в неискаженном пространстве.

\subsubsection{Практическая реализация}

Рабочий процесс состоит из двух этапов:
\begin{enumerate}
    \item \textbf{Калибровка}: на ПЗС-линейку проецируется образ зебры; по считанному сигналу определяется $k^*$ и заносится в модель.

    \begin{figure}[ht!]
        \centering
        \includegraphics[width=0.5\linewidth]{BespyatyyIV/images/Зебра.jpg}
        \caption{Зебра}
        \label{fig:zebra}
    \end{figure}
    
    Следует отметить, что на практике оценка коэффициента $k^*$ выполнялась не по всем обнаруженным фронтам, а только по одноименным, то есть либо только по восходящим, либо только по нисходящим (см. рис. \ref{fig:distortion_calibration}). Такой подход связан с тем, что при изготовлении растра методом 3D-печати размеры пропусков и заполнений не могут быть воспроизведены строго одинаковыми, поэтому расстояния между разноименными фронтами содержат дополнительную систематическую составляющую, не связанную с дисторсией оптической системы.

    \begin{figure}[ht!]
        \centering
        \includegraphics[width=\linewidth]{BespyatyyIV/images/дисторсия_калибровка.pdf}
        \caption{Диагностические выводы процесса калибровки}
        \label{fig:distortion_calibration}
    \end{figure}

    Оценка коэффициента дисторсии, проведенная в процессе калибровки серией из трех измерений, для первой головки была принята равной $k_1^* \simeq 0.052$, для второй $k_2^* \simeq 0.050$

    \item \textbf{Рабочий цикл}: при считывании произвольного растра значение $k^*$ используется для геометрической коррекции сигнала перед этапом битовой разметки (см. рис \ref{fig:distortion_correction}).

    \begin{figure}[ht!]
        \centering
        \includegraphics[width=\linewidth]{BespyatyyIV/images/пример_восстановления_дисторсии.png}
        \caption{Пример восстановления дисторсии в сигнале с матрицы после калибровки при $k = 0.05$}
        \label{fig:distortion_correction}
    \end{figure}
\end{enumerate}

\subsection{Модель формирования, обработки и восстановления сигнала}\label{subsec:model}

Для исследования процесса считывания кодового растра и последующего восстановления битовой последовательности рассматривается единая модель формирования и обработки сигнала линейной ПЗС-матрицы. В рамках этой модели используются два типа входных данных: синтетический сигнал, формируемый численно по заданной битовой последовательности, и экспериментальный сигнал, непрерывно считываемый с матрицы в реальной оптической схеме. Различие между ними состоит в способе получения исходного сигнала, тогда как дальнейшая обработка, выделение фронтов и восстановление битовой структуры выполняются по единому алгоритму.

В синтетическом случае модель позволяет формировать сигнал с заранее известной разметкой, что необходимо для отладки параметров обработки и оценки точности восстановления. В экспериментальном случае исходным является реальный сигнал с линейной ПЗС-матрицы.

\subsubsection{Геометрическая коррекция дисторсии}

В соответствии с рассмотренной ранее моделью дисторсии (см. подраздел~\ref{subsec:distortion_model}), перед дальнейшей обработкой экспериментального сигнала выполняется геометрическая коррекция, переводящая его в неискаженное координатное пространство. Для синтетического сигнала этот этап либо не требуется, либо учитывается уже на стадии формирования. Далее в обоих случаях рассматривается единый сигнал $\tilde S$, заданный на равномерной пиксельной сетке и подготовленный к последующей обработке.


\subsubsection{Нормализация и сглаживание}

\begin{enumerate}
    \item \textbf{Глобальная нормализация.} Исправленный сигнал $\tilde S$ приводится к единичному диапазону:
    \begin{equation}
        Q = \frac{\tilde S - \min \tilde S}{\max \tilde S - \min \tilde S + \varepsilon}.
    \end{equation}

    \item \textbf{Гауссово сглаживание.} К нормированному сигналу применяется гауссово сглаживание с параметром $\sigma_3$, подавляющее высокочастотный шум без существенного смещения положений фронтов:
    \begin{equation}
        Q' = G(\sigma) * Q.
    \end{equation}

    \item \textbf{Коррекция виньетирования.} Сглаженный сигнал делится на нормированный предварительно измеренный профиль виньетирования, после чего вновь выполняется глобальная нормализация:
    \begin{equation}
        \hat Q = \mathrm{norm}\!\left(\frac{Q'}{\tilde V}\right),
        \quad
        \tilde V = \mathrm{norm}(V).
    \end{equation}
    Данная операция компенсирует неравномерность амплитуды сигнала вдоль матрицы, обусловленную виньетированием объектива, и тем самым приводит сигнал к виду, более удобному для последующего выделения фронтов.
\end{enumerate}

\subsubsection{Фильтр Блэ--Риу и вычисление производных}

Для оценки локальных изменений сигнала применяется КИХ-фильтр Блэ--Риу~\cite{blais_rioux_peak}. Ядро фильтра порядка $N$ имеет длину $2N+1$:
\begin{equation}
    \mathrm{BR} = \frac{1}{2N}\bigl[\underbrace{-1,\ldots,-1}_{N},\;0,\;\underbrace{1,\ldots,1}_{N}\bigr].
\end{equation}
Свертка этого ядра с сигналом $\hat Q$ дает сглаженную первую производную:
\begin{equation}
    D_1 = \mathrm{BR} * \hat Q.
\end{equation}
Повторное применение того же ядра к $D_1$ дает сглаженную вторую производную:
\begin{equation}
    D_2 = \mathrm{BR} * D_1.
\end{equation}

Экстремумы первой производной соответствуют наиболее резким переходам между соседними участками растра. Вторая производная используется как вспомогательная величина для точного определения положений этих переходов.

\subsubsection{Обнаружение фронтов}

\begin{enumerate}
    \item \textbf{Определение нулей второй производной.} Положения фронтов соответствуют экстремумам $D_1$, то есть нулям $D_2$. Эти положения определяются по точкам смены знака с последующим субпиксельным уточнением линейной интерполяцией:
    \begin{equation}
        j:\;D_2[j]\cdot D_2[j+1] < 0
        \;\Rightarrow\;
        p_i \simeq j + \frac{D_2[j]}{D_2[j]-D_2[j+1]}.
        \label{eq:zerocross}
    \end{equation}
    Значение $D_1$ в точке $p_i$ вычисляется той же линейной интерполяцией между $D_1[j]$ и $D_1[j+1]$.

    \item \textbf{Пороговая фильтрация по амплитуде.} Из числа найденных кандидатов сохраняются только те, для которых абсолютное значение первой производной превышает долю $\alpha$ от максимального значения по всем кандидатам:
    \begin{equation}
        |D_1(p_i)| \geq \alpha \cdot \max_i |D_1(p_i)|.
    \end{equation}

    \item \textbf{Фильтрация по минимальному расстоянию.} Среди кандидатов, расположенных друг к другу ближе чем на $\delta = \beta\cdot k$ пикселей, где $k$ --- ширина бита в пикселях, а $\beta$ --- настраиваемый коэффициент, сохраняется кандидат с максимальным значением $|D_1|$:
    \begin{equation}
        |p_i - p_j| < \delta
        \;\Rightarrow\;
        \text{оставить}\;\arg\max |D_1|.
    \end{equation}
\end{enumerate}

Результатом этапа является упорядоченный массив положений фронтов и соответствующих значений первой производной. Тем самым формируется геометрическая разметка структуры сигнала, отражающая положения переходов между соседними элементами кодового растра.

\subsubsection{Восстановление битовой последовательности}

\begin{enumerate}
    \item \textbf{Построение битовых сегментов.} По массиву обнаруженных фронтов $\{p_i\}$ строятся сегменты $[p_i,\,p_{i+1})$. Для каждого сегмента число битов определяется как ближайшее целое к отношению его длины к априорной ширине бита $k$:
    \begin{equation}
        n_i = \left[\frac{p_{i+1}-p_i}{k}\right].
    \end{equation}
    Значение бита внутри сегмента определяется знаком $D_1$ в точке левого фронта:
    \begin{equation}
        \mathrm{val}_i = \begin{cases}
            1, & D_1(p_i) > 0,\\
            0, & D_1(p_i) < 0.
        \end{cases}
    \end{equation}

    \item \textbf{Оценка локального периода.} Для каждого видимого сегмента определяется локальный период
    \begin{equation}
        k_i = \frac{p_{i+1}-p_i}{n_i},
    \end{equation}

    \item \textbf{Восстановление положений битов.} Внутри каждого сегмента левые границы битов распределяются равномерно с периодом $k_i$:
    \begin{equation}
        p_i^{\,j} = p_i + j\cdot k_i,
        \quad 0 \leq j < n_i.
    \end{equation}

    \item \textbf{Оценка точности восстановления битовой последовательности\footnote{\label{ftnt:model}В цифровой модели}.} Для оценки достоверности восстановленной последовательности применяется скользящее окно сравнения: восстановленный массив битов сдвигается относительно эталонного, и для каждого смещения подсчитывается число совпадений. Наилучшее смещение определяет оценку точности:
    \begin{equation}
        \mathrm{Acc} = \frac{\max_{\Delta}\sum_j \mathbf{1}[\hat b_{j+\Delta} = b_j]}{n_{\min}}, \quad n_{\min} = \min(|\hat B|,\,|B|)
    \end{equation}

    \item \textbf{Оценка погрешности положений фронтов.}\textsuperscript{\ref{ftnt:model}} Для каждого истинного фронта $t_i$ из эталонного массива методом ближайшего соседа, при ограничении $|p - t_i| < k/2$, выбирается соответствующий обнаруженный фронт, после чего вычисляется ошибка
    \begin{equation}
        e_i = p_{\mathrm{nearest}} - t_i.
    \end{equation}
    По массиву ошибок $\{e_i\}$ определяется среднеквадратическая погрешность:
    \begin{equation}
        \sigma_e = \sqrt{\frac{1}{|\{e_i\}|}\sum_i e_i^2}.
    \end{equation}
\end{enumerate}

\subsubsection{Результаты на синтетическом сигнале}
\begin{figure}[ht!]
    \centering
    \includegraphics[width=\linewidth]{BespyatyyIV/images/Цифровая_модель.pdf}
    \caption{Полный вывод с диагностическими графиками в синтетической модели. Было достигнуто СКО определения положений фронтов менее $0.1$ пиксела, что составляет менее $2\%$ ширины бита. Точность восстановления битовой последовательности составила $100\%$ при отношении сигнал/шум не менее $10$ дБ.}
    \label{fig:ыynth_model}
\end{figure}

В ходе статистического анализа работы модели на синтетическом сигнале в качестве рабочего был выбран порядок алгоритма Блэ-Риу равным $N=2$. Диагностические графики приведены ниже (рис.\ref{fig:wplots} и \ref{fig:bnplots}):

\begin{figure}[ht!]
    \centering
    \begin{subfigure}[b]{0.48\textwidth}
        \centering
        \includegraphics[width=\textwidth]{BespyatyyIV/images/diagnostic_blur=0.05_noise=200.pdf}
        \caption{Уровень шума ок. 13 дБ, уровень размытия 5 процентов от ширины бита}
        \label{fig:b0.05n200}
    \end{subfigure}
    \hfill
    \begin{subfigure}[b]{0.48\textwidth}
        \centering
        \includegraphics[width=\textwidth]{BespyatyyIV/images/diagnostic_blur=0.2_noise=200.pdf}
        \caption{Уровень шума ок. 13 дБ, уровень размытия 20 процентов от ширины бита}
        \label{fig:b0.2n200}
    \end{subfigure}

    \caption{Оцененная зависимость СКО алгоритма Блэ-Риу с $N=2$ и $N=4$ от ширины бита в пикселах}
    \label{fig:wplots}
\end{figure}

\begin{figure}[ht!]
    \centering
    \begin{subfigure}[b]{0.48\textwidth}
        \centering
        \includegraphics[width=\textwidth]{BespyatyyIV/images/diagnostic_N=2_bit=6.pdf}
        \caption{$N=2$}
        \label{fig:n2}
    \end{subfigure}
    \hfill
    \begin{subfigure}[b]{0.48\textwidth}
        \centering
        \includegraphics[width=\textwidth]{BespyatyyIV/images/diagnostic_N=4_bit=6.pdf}
        \caption{$N=4$}
        \label{fig:bnplots}
    \end{subfigure}

    \caption{Тепловая карта СКО от уровня шума (до 12 дБ) и уровня размытия перед обработкой в относительных единицах ширины бита}
    \label{fig:two-sub}
\end{figure}

\subsubsection{Результаты на действительном сигнале}
\begin{figure}[ht!]
    \centering
    \includegraphics[width=\linewidth]{BespyatyyIV/images/Blais-Rioux Edge Detector.png}
    \caption{Пример работы и восстановления на сигнале с матрицы}
    \label{fig:phys_model}
\end{figure}

\subsection{Оценка абсолютного угла диска}

После восстановления битовой последовательности по сигналу ПЗС"=матрицы выполняется оценка абсолютного угла поворота кодового диска энкодера. На вход алгоритма подаются результаты разметки, полученные на предыдущем этапе обработки сигнала \ref{subsec:model}.

\subsubsection{Кодовая дорожка и таблица кодовых слов}

Кодовая дорожка диска представляет собой двоичную последовательность де Брейна длиной $N_{\mathrm{disk}}$ бит. Подробное описание см. в подразделе \ref{subsubsec:code_structure}.

Для каждой начальной позиции $s \in [0,\,N_{\mathrm{disk}})$ из кодовой последовательности извлекается слово длиной $m$ с циклическим обходом, после чего ему ставится в соответствие угловое положение центра данного окна:
\begin{equation}
    \varphi_s = \left(s + \frac{m}{2}\right)\cdot\frac{360^\circ}{N_{\mathrm{disk}}}.
    \label{eq:cw_angle}
\end{equation}
Поскольку используемая последовательность является последовательностью де Брейна, все $N_{\mathrm{disk}}$ кодовых слов различны, и построенная таблица соответствует взаимно однозначному отображению между кодовым словом и угловой координатой на диске.

\subsubsection{Масштабный коэффициент угол/пиксель}

Для перехода от пиксельных координат к угловым используется масштабный коэффициент $\lambda$, определяемый через средний период бита $\bar k$ в пикселях. Эта величина оценивается по битовым сегментам, выделенным на этапе восстановления:
\begin{equation}
    \bar k = \frac{\sum_i n_i\,k_i}{\sum_i n_i}, \quad k_i = \frac{l_i}{n_i},
\end{equation}
где $n_i$ --- число битов в $i$"=м сегменте, $l_i$ --- длина сегмента в пикселях, а $k_i$ --- локальная оценка периода бита в данном сегменте. Тогда угловой масштаб определяется выражением
\begin{equation}
    \lambda = \frac{360^\circ}{N_{\mathrm{disk}}\cdot\bar k}.
    \label{eq:angle_per_px}
\end{equation}

\subsubsection{Скользящее окно и оценка угла}

Для последующей оценки угла используются только те восстановленные биты, которые целиком принадлежат соответствующему сегменту с учетом малого допуска $\varepsilon$, связанного с конечной точностью субпиксельной разметки:
\begin{equation}
    p_{i,\,\mathrm{start}} - \varepsilon \leq b_{\mathrm{start}} \leq b_{\mathrm{end}} \leq p_{i,\,\mathrm{end}} + \varepsilon.
\end{equation}
Для каждого корректно определенного бита сохраняются его бинарное значение $b_j \in \{0,1\}$, координата центра $x_j = b_{\mathrm{start}} + k_j/2$ и ширина $k_j$, равная локальному периоду соответствующего сегмента.

По последовательности валидных битов строятся скользящие окна длиной $m$ бит. Если число валидных битов равно $n$, то число таких окон составляет $n_w = n - m + 1$. Центр $j$"=го окна определяется как взвешенное среднее координат входящих в него битов, где весом служит ширина соответствующего бита:
\begin{equation}
    x_{\mathrm{win},j} = \frac{\sum_{i=j}^{j+m-1} k_i\,x_i}{\sum_{i=j}^{j+m-1} k_i}.
    \label{eq:win_center}
\end{equation}

Для каждого окна формируется кодовое слово $c_j$, составленное из значений входящих в него битов, и выполняется поиск этого слова в предварительно построенной таблице кодовых слов. Если найдено соответствие позиции $s$ на диске, то табличное угловое положение окна определяется по формуле~\eqref{eq:cw_angle}.

Угол центра поля зрения матрицы равен углу центра распознанного кодового слова на диске за вычетом углового смещения этого слова относительно центра матрицы. Таким образом, оценка абсолютного угла диска по данному окну вычисляется по смещению центра окна относительно центра матрицы $x_0$:
\begin{equation*}
    \theta_j = \bigl(\varphi_j - \lambda\,(x_{\mathrm{win},j} - x_0)\bigr) \bmod 360^\circ.
\end{equation*}

\subsubsection{Весовая функция виньетирования}

Биты, расположенные вблизи краев поля зрения матрицы значительнее всего были подвержены обработке как по коррекции дисторсии, так по виньетированию. По этой причине кодовые окна, центры которых находятся ближе к краям матрицы, определены с потенциально большей ошибкой, чем центральные. Поэтому в итоговую оценку угла им следует учитываться с различным весом.


где  Вес окна задается модельной функцией виньетирования:
\begin{equation}
    w_j = \begin{cases}
        \max\!\left(0,\;1 - \kappa\!\left(1 - \left|\cos\!\dfrac{\pi u_j}{2}\right|^s\right)\right), & |u_j| \leq 1,\\[6pt]
        0, & |u_j| > 1,
    \end{cases}, \quad  u_j = \frac{x_{\mathrm{win},j} - x_0}{L/2}
    \label{eq:vign_weight}
\end{equation}
где $\kappa$ характеризует глубину виньетирования, $L$~--- ширина матрицы в пикселях, $s$ определяет форму убывания веса к краям поля зрения. При $\kappa = 0$ веса всех окон совпадают и равны единице.

\subsubsection{Фильтрация выбросов}

Среди оценок $\theta_j$ могут присутствовать выбросы, обусловленные ошибками восстановления отдельных битов. Для их исключения используется итерационная процедура, учитывающая циклическую структуру угловых величин, что критично в окрестности границы $0^\circ/360^\circ$:

\begin{enumerate}
    \item по текущему набору точек вычисляются взвешенное круговое среднее $\bar\theta$ и круговое стандартное отклонение $\sigma_\theta$;
    \item для каждой оценки определяется минимальное расстояние до среднего:
    \begin{equation}
        d_j = \min\nolimits_{k\in\mathbb{Z}} |\theta_j - \bar\theta + k\cdot 360^\circ|;
    \end{equation}
    \item исключаются оценки, для которых выполняется условие
    \begin{equation}
        d_j > \max(\sigma_\theta,\,\sigma_{\min}),
    \end{equation}
    где $\sigma_{\min}$ --- минимально допустимое значение порога, предотвращающее чрезмерное отбрасывание оценок при малом разбросе;
    \item шаги 1--3 повторяются до стабилизации набора точек либо до достижения максимального числа итераций.
\end{enumerate}

\subsubsection{Итоговая оценка угла}

По отфильтрованному набору оценок $\{\theta_j\}$ с весами $\{w_j\}$ вычисляется итоговая оценка абсолютного угла диска как взвешенное круговое среднее:
\begin{equation}
    \bar\theta = \mathrm{arg}\!\left(\sum_j w_j \, e^{\,i\,\theta_j\cdot 2\pi/360^\circ}\right).
    \label{eq:circ_mean}
\end{equation}

Мерой согласованности полученных оценок служит круговое стандартное отклонение, определяемое через нормированную длину результирующего вектора $R$:
\begin{equation}
    R = \frac{\left|\sum_j w_j \, e^{\,i\,\theta_j\cdot 2\pi/360^\circ}\right|}{\sum_j w_j}, \qquad \sigma_\theta = \sqrt{-2\ln R}
    \label{eq:circ_std}
\end{equation}